

\section{Create a new dataset: {\tt streamOpenWrite}}
\index{streamOpenWrite}
\label{streamOpenWrite}

The function {\tt streamOpenWrite} creates a new datset.
\subsection*{Usage}

\begin{verbatim}
    INTEGER FUNCTION streamOpenWrite(CHARACTER*(*) path, INTEGER filetype)
\end{verbatim}

\hspace*{4mm}\begin{minipage}[]{15cm}
\begin{deflist}{\tt filetype\ }
\item[{\tt path}]
The name of the new dataset
\item[{\tt filetype}]
The type of the file format, one of the set of predefined CDI file format types.
                     The valid CDI file format types are {\tt FILETYPE\_GRB}, {\tt FILETYPE\_NC},
                     {\tt FILETYPE\_NC2}, {\tt FILETYPE\_SRV}, {\tt FILETYPE\_EXT} and {\tt FILETYPE\_IEG}.

\end{deflist}
\end{minipage}

\subsection*{Result}

Upon successful completion {\tt streamOpenWrite} returns an identifier to the
open stream. Otherwise, a negative number with the error status is returned.


\subsection*{Errors}


\hspace*{4mm}\begin{minipage}[]{15cm}
\begin{deflist}{\tt CDI\_EUFILETYPE\ }
\item[{\tt CDI\_ESYSTEM}]
Operating system error
\item[{\tt CDI\_EINVAL}]
Invalid argument
\item[{\tt CDI\_EUFILETYPE}]
Unsupported file type
\item[{\tt CDI\_ELIBNAVAIL}]
Library not available
\end{deflist}
\end{minipage}


\subsection*{Example}

Here is an example using {\tt streamOpenWrite} to create a new netCDF file 
named {\tt foo.nc} for writing:

\begin{lstlisting}[backgroundcolor=\color{zebg}, basicstyle=\footnotesize]
    INCLUDE 'cdi.h'
       ...
    INTEGER streamID
       ...
    streamID = streamOpenWrite("foo.nc", FILETYPE_NC)
    IF ( streamID .LT. 0 ) CALL handle_error(streamID)
       ...
\end{lstlisting}


\section{Open a dataset for reading: {\tt streamOpenRead}}
\index{streamOpenRead}
\label{streamOpenRead}

The function {\tt streamOpenRead} opens an existing dataset for reading.

\subsection*{Usage}

\begin{verbatim}
    INTEGER FUNCTION streamOpenRead(CHARACTER*(*) path)
\end{verbatim}

\hspace*{4mm}\begin{minipage}[]{15cm}
\begin{deflist}{\tt path\ }
\item[{\tt path}]
The name of the dataset to be read

\end{deflist}
\end{minipage}

\subsection*{Result}

Upon successful completion {\tt streamOpenRead} returns an identifier to the
open stream. Otherwise, a negative number with the error status is returned.


\subsection*{Errors}


\hspace*{4mm}\begin{minipage}[]{15cm}
\begin{deflist}{\tt CDI\_EUFILETYPE\ }
\item[{\tt CDI\_ESYSTEM}]
Operating system error
\item[{\tt CDI\_EINVAL}]
Invalid argument
\item[{\tt CDI\_EUFILETYPE}]
Unsupported file type
\item[{\tt CDI\_ELIBNAVAIL}]
Library not available
\end{deflist}
\end{minipage}


\subsection*{Example}

Here is an example using {\tt streamOpenRead} to open an existing netCDF
file named {\tt foo.nc} for reading:

\begin{lstlisting}[backgroundcolor=\color{zebg}, basicstyle=\footnotesize]
    INCLUDE 'cdi.h'
       ...
    INTEGER streamID
       ...
    streamID = streamOpenRead("foo.nc")
    IF ( streamID .LT. 0 ) CALL handle_error(streamID)
       ...
\end{lstlisting}


\section{Close an open dataset: {\tt streamClose}}
\index{streamClose}
\label{streamClose}

The function {\tt streamClose} closes an open dataset.

\subsection*{Usage}

\begin{verbatim}
    SUBROUTINE streamClose(INTEGER streamID)
\end{verbatim}

\hspace*{4mm}\begin{minipage}[]{15cm}
\begin{deflist}{\tt streamID\ }
\item[{\tt streamID}]
Identifier, from a previous call to {\tt streamOpenRead}
                     or {\tt streamOpenWrite}

\end{deflist}
\end{minipage}


\section{Get the filetype: {\tt streamInqFiletype}}
\index{streamInqFiletype}
\label{streamInqFiletype}

The function {\tt streamInqFiletype} returns the filetype of a stream.

\subsection*{Usage}

\begin{verbatim}
    INTEGER FUNCTION streamInqFiletype(INTEGER streamID)
\end{verbatim}

\hspace*{4mm}\begin{minipage}[]{15cm}
\begin{deflist}{\tt streamID\ }
\item[{\tt streamID}]
Stream identifier

\end{deflist}
\end{minipage}

\subsection*{Result}

{\tt streamInqFiletype} returns the type of the file format,
one of the set of predefined CDI file format types.
The valid CDI file format types are {\tt FILETYPE\_GRB}, {\tt FILETYPE\_NC},
{\tt FILETYPE\_NC2}, {\tt FILETYPE\_SRV}, {\tt FILETYPE\_EXT} and {\tt FILETYPE\_IEG}.



\section{Define the variable list: {\tt streamDefVlist}}
\index{streamDefVlist}
\label{streamDefVlist}

The function {\tt streamDefVlist} defines the variable list of a stream.

\subsection*{Usage}

\begin{verbatim}
    SUBROUTINE streamDefVlist(INTEGER streamID, INTEGER vlistID)
\end{verbatim}

\hspace*{4mm}\begin{minipage}[]{15cm}
\begin{deflist}{\tt streamID\ }
\item[{\tt streamID}]
Stream identifier
\item[{\tt vlistID}]
Variable list identifier

\end{deflist}
\end{minipage}


\section{Get the variable list: {\tt streamInqVlist}}
\index{streamInqVlist}
\label{streamInqVlist}

The function {\tt streamInqVlist} returns the variable list of a stream.

\subsection*{Usage}

\begin{verbatim}
    INTEGER FUNCTION streamInqVlist(INTEGER streamID)
\end{verbatim}

\hspace*{4mm}\begin{minipage}[]{15cm}
\begin{deflist}{\tt streamID\ }
\item[{\tt streamID}]
Stream identifier

\end{deflist}
\end{minipage}

\subsection*{Result}

{\tt streamInqVlist} returns an identifier to the variable list.



\section{Define time step: {\tt streamDefTimestep}}
\index{streamDefTimestep}
\label{streamDefTimestep}

The function {\tt streamDefTimestep} defines the time step of a stream.

\subsection*{Usage}

\begin{verbatim}
    INTEGER FUNCTION streamDefTimestep(INTEGER streamID, INTEGER tsID)
\end{verbatim}

\hspace*{4mm}\begin{minipage}[]{15cm}
\begin{deflist}{\tt streamID\ }
\item[{\tt streamID}]
Stream identifier
\item[{\tt tsID}]
Timestep identifier

\end{deflist}
\end{minipage}

\subsection*{Result}

{\tt streamDefTimestep} returns the number of records of the time step.



\section{Get time step: {\tt streamInqTimestep}}
\index{streamInqTimestep}
\label{streamInqTimestep}

The function {\tt streamInqTimestep} returns the time step of a stream.

\subsection*{Usage}

\begin{verbatim}
    INTEGER FUNCTION streamInqTimestep(INTEGER streamID, INTEGER tsID)
\end{verbatim}

\hspace*{4mm}\begin{minipage}[]{15cm}
\begin{deflist}{\tt streamID\ }
\item[{\tt streamID}]
Stream identifier
\item[{\tt tsID}]
Timestep identifier

\end{deflist}
\end{minipage}

\subsection*{Result}

{\tt streamInqTimestep} returns the number of records of the time step.



\section{Write a variable: {\tt streamWriteVar}}
\index{streamWriteVar}
\label{streamWriteVar}

The function streamWriteVar writes the values of one time step of a variable 
to an open dataset.
\subsection*{Usage}

\begin{verbatim}
    SUBROUTINE streamWriteVar(INTEGER streamID, INTEGER varID, REAL*8 data, 
                              INTEGER nmiss)
\end{verbatim}

\hspace*{4mm}\begin{minipage}[]{15cm}
\begin{deflist}{\tt streamID\ }
\item[{\tt streamID}]
Stream identifier, from a previous call to streamOpenWrite
\item[{\tt varID}]
Variable identifier
\item[{\tt data}]
Pointer to a block of data values to be written
\item[{\tt nmiss}]
Number of missing values

\end{deflist}
\end{minipage}


\section{Read a variable: {\tt streamReadVar}}
\index{streamReadVar}
\label{streamReadVar}

The function streamReadVar reads all the values of one time step of a variable
from an open dataset.
\subsection*{Usage}

\begin{verbatim}
    SUBROUTINE streamReadVar(INTEGER streamID, INTEGER varID, REAL*8 data, 
                             INTEGER nmiss)
\end{verbatim}

\hspace*{4mm}\begin{minipage}[]{15cm}
\begin{deflist}{\tt streamID\ }
\item[{\tt streamID}]
Stream identifier, from a previous call to streamOpenRead
\item[{\tt varID}]
Variable identifier
\item[{\tt data}]
Pointer to the location into which the data value is read
\item[{\tt nmiss}]
Number of missing values

\end{deflist}
\end{minipage}


\section{Write a horizontal slice of a variable: {\tt streamWriteVarSlice}}
\index{streamWriteVarSlice}
\label{streamWriteVarSlice}

The function streamWriteVarSlice writes the values of a horizontal slice of a 
variable to an open dataset.
\subsection*{Usage}

\begin{verbatim}
    SUBROUTINE streamWriteVarSlice(INTEGER streamID, INTEGER varID, INTEGER levelID, 
                                   REAL*8 data, INTEGER nmiss)
\end{verbatim}

\hspace*{4mm}\begin{minipage}[]{15cm}
\begin{deflist}{\tt streamID\ }
\item[{\tt streamID}]
Stream identifier, from a previous call to streamOpenWrite
\item[{\tt varID}]
Variable identifier
\item[{\tt levelID}]
Level identifier
\item[{\tt data}]
Pointer to a block of data values to be written
\item[{\tt nmiss}]
Number of missing values

\end{deflist}
\end{minipage}


\section{Read a horizontal slice of a variable: {\tt streamReadVarSlice}}
\index{streamReadVarSlice}
\label{streamReadVarSlice}

The function streamReadVar reads all the values of a horizontal slice of a variable
from an open dataset.
\subsection*{Usage}

\begin{verbatim}
    SUBROUTINE streamReadVarSlice(INTEGER streamID, INTEGER varID, INTEGER levelID, 
                                  REAL*8 data, INTEGER nmiss)
\end{verbatim}

\hspace*{4mm}\begin{minipage}[]{15cm}
\begin{deflist}{\tt streamID\ }
\item[{\tt streamID}]
Stream identifier, from a previous call to streamOpenRead
\item[{\tt varID}]
Variable identifier
\item[{\tt levelID}]
Level identifier
\item[{\tt data}]
Pointer to the location into which the data value is read
\item[{\tt nmiss}]
Number of missing values

\end{deflist}
\end{minipage}
